Dues mostres radioactives tenen, en un moment donat, $1,00\times 10^{-1}$~mol cadascuna. Les mostres són de dos isòtops diferents de l'element radó (Rn): en concret, de radó 222 ($^{222}$Rn) i de radó 224 ($^{224}$Rn). Els dos isòtops són radioactius i tenen, respectivament, períodes de semidesintegració de 3,82 dies i 1,80 hores. El primer presenta una desintegració de tipus $\alpha$ i el nucli fill és un isòtop del poloni (Po), mentre que el segon presenta una desintegració de tipus $\beta^-$ i el nucli fill és un isòtop del franci (Fr).

\begin{apartat}{1}
Escriviu les equacions nuclears de les dues desintegracions radioactives amb totes les partícules que hi intervenen i els seus nombres atòmics i màssics. Calculeu quants àtoms de $^{224}$Rn no s'hauran desintegrat encara quan restin $9,00\times 10^{-2}$~mol de la mostra del $^{222}$Rn per desintegrar-se.
\end{apartat}

\begin{apartat}{1}
L'energia que es desprèn per cada desintegració d'un nucli de $^{222}$Rn és de 5,590~MeV. Calculeu el defecte de massa d'aquesta reacció nuclear.
\end{apartat}

\dades{%
  Nombre d'Avogadro, $N_\mathrm{A} = 6,022\times 10^{23}\ \mathrm{mol^{-1}}$.\\
  Velocitat de la llum, $c = 3,00\times 10^{8}\ \mathrm{m\,s^{-1}}$.\\
  $1\ \mathrm{eV} = 1,60\times 10^{-19}\ \mathrm{J}$.\\
  Nombre atòmic del radó $= 86$.}
